% Client session phases of Qadra and their effect on in-memory drafts.
% Phases and transitions follow web/apps/qadra/src/lib/session-lifecycle.mjs;
% draft effects follow web/apps/qadra/src/lib/draft-registry.mjs.
% Layout: side columns are 4.2 cm wide so every title fits on one line,
% the middle column is 3.8 cm wide, and the figure spans 16.7 cm.
\begin{tikzpicture}[
  phase/.style={
    draw=ipnvino, rounded corners=1.5mm, line width=0.6pt,
    text width=4.2cm, align=center, inner sep=4pt,
    minimum height=2.3cm, font=\footnotesize,
    execute at begin node={\hyphenpenalty=10000\exhyphenpenalty=10000}
  },
  middle/.style={text width=3.8cm},
  accepted/.style={minimum height=3.3cm},
  event/.style={
    phase, middle, draw=black!45, fill=black!3
  },
  note/.style={
    draw=black!45, dashed, rounded corners=1.5mm, fill=black!3,
    text width=7.9cm, align=flush left, inner sep=4pt, font=\footnotesize,
    execute at begin node={\hyphenpenalty=10000\exhyphenpenalty=10000}
  },
  flow/.style={-{Stealth[length=2mm]}, line width=0.6pt},
  lab/.style={font=\footnotesize, align=center, inner sep=2pt}
]
  % Row 1: no session, authentication and expired session.
  \node[phase] (out) at (-5.95,0) {
    \textbf{Sin sesión} (\texttt{signed-out})\\[2pt]
    Guarda cerrada;\\sin credencial\\[2pt]
    Borradores: ninguno
  };
  \node[event] (mfa) at (0,0) {
    \textbf{Contraseña o certificado}\\
    \textbf{Owner, y MFA}\\[2pt]
    Plazos confirmados\\con reloj monótono\\[2pt]
    Misma cuenta:\\conserva las capturas\\
    Otra cuenta: las descarta
  };
  \node[phase] (exp) at (5.95,0) {
    \textbf{Vencida} (\texttt{expired})\\[2pt]
    Guarda cerrada;\\credencial retirada\\[2pt]
    Borradores: captura síncrona\\de los editores registrados
  };
  % Row 2: accepted session.
  \node[phase, accepted, fill=ipnvino!6] (act) at (-5.95,-4.65) {
    \textbf{Activa} (\texttt{active})\\[2pt]
    Guarda abierta hasta\\el plazo confirmado\\[2pt]
    Borradores: recuperación\\a petición, tras consultar\\de nuevo el destino
  };
  \node[phase, accepted, middle] (hid) at (0,-4.65) {
    \textbf{Oculta} (\texttt{hidden})\\[2pt]
    Guarda cerrada:\\admisión suspendida\\[2pt]
    Borradores: sin cambios
  };
  \node[phase, accepted] (chk) at (5.95,-4.65) {
    \textbf{Comprobación} (\texttt{checking})\\[2pt]
    Guarda cerrada;\\consulta nueva del estado\\[2pt]
    Error de red: cerrada\\hasta reintento o cierre\\[2pt]
    Borradores: sin cambios
  };
  % Accepted-session frame.
  \draw[dashed, black!45, rounded corners=2mm]
    (-8.35,-2.8) rectangle (8.35,-7.3);
  \node[lab, anchor=south west, text=black!70] at (-8.3,-7.27)
    {Sesión aceptada};

  % Authentication.
  \draw[flow] (out.east) -- node[lab, above] {inicio} (mfa.west);
  \draw[flow] (exp.west) -- node[lab, above] {reingreso} (mfa.east);
  \draw[flow] (mfa.south west) -- node[lab, above left]
    {pestaña\\visible} (act.north east);
  \draw[flow] (mfa.south) -- node[lab, right] {pestaña oculta} (hid.north);

  % Visibility cycle.
  \draw[flow] (act.east) -- node[lab, above] {pestaña\\oculta} (hid.west);
  \draw[flow] ([yshift=3mm]hid.east) -- node[lab, above] {pestaña\\visible}
    ([yshift=3mm]chk.west);
  \draw[flow] ([yshift=-3mm]chk.west) -- node[lab, below] {pestaña\\oculta}
    ([yshift=-3mm]hid.east);
  \draw[flow] (chk.south) -- ++(0,-0.4) -| node[lab, pos=0.25, below]
    {estado confirmado} ([xshift=6mm]act.south);

  % Exits from any accepted phase, from the frame up to the phase box.
  \coordinate (exitl) at (-7.95,-2.8);
  \coordinate (exitr) at (7.95,-2.8);
  \draw[flow] (exitl) -- node[lab, right] {cierre voluntario:\\descarta borradores}
    (exitl |- out.south);
  \draw[flow] (exitr) -- node[lab, left] {plazo vencido, rechazo\\o estado inválido}
    (exitr |- exp.south);

  % Notes.
  \node[note, anchor=north west] at (-8.35,-7.55) {
    Guarda cerrada: Qadra rechaza localmente las solicitudes
    protegidas; sólo la consulta de estado y la revocación de
    una credencial vigente quedan exentas.
  };
  \node[note, anchor=north east] at (8.35,-7.55) {
    Desmontada (\texttt{disposed}): al desmontarse la aplicación,
    desde cualquier fase, descarta los borradores y retira la
    credencial; la admisión queda cerrada.
  };
\end{tikzpicture}
